% Part: computability
% Chapter: computability-theory
% Section: ce-closed-cup-cap

\documentclass[../../../include/open-logic-section]{subfiles}

\begin{document}

% SEGMENT OLP-0241-S01

\olfileid{cmp}{thy}{clo}

\olsection[க. எ. கணங்களின் ஒன்றிப்பும் வெட்டும்]{கணிக்கத்தக்கவாறு
  எண்ணிடத்தக்க கணங்கள் ஒன்றிப்பின்கீழும் வெட்டின்கீழும் மூடியவை}

கணிக்கத்தக்கவாறு எண்ணிடத்தக்க கணங்களின் சில மூடுகைப் பண்புகளைப்
பின்வரும் தேற்றம் வழங்குகிறது.

\begin{thm}
$A$ உம் $B$ உம் கணிக்கத்தக்கவாறு எண்ணிடத்தக்கவை எனக் கொள்க.
அப்போது $A \cap B$ உம் $A \cup B$ உம் அவ்வாறே உள்ளன.
\end{thm}

\begin{proof}
கணிக்கத்தக்கவாறு எண்ணிடத்தக்க கணங்களின் பல்வேறு இயல்பாக்கங்களைப்
பயன்படுத்த \olref[eqc]{thm:ce-equiv} அனுமதிக்கிறது. விளக்கத்திற்காக,
சில வெவ்வேறு நிரூபணங்களை வழங்குவோம்.

முதல் நிரூபணத்திற்கு, $A$ ஆனது கணிக்கத்தக்க சார்பு~$f$ ஆலும்
$B$ ஆனது கணிக்கத்தக்க சார்பு~$g$ ஆலும் எண்ணிடப்படுகின்றன எனக் கொள்க.
பின்வருமாறு எடுத்துக்கொள்க:
\begin{align*}
h(x) & = \umin{y}{(f(y) = x \lor g(y) = x)} \text{ மற்றும்}\\
j(x) & = \umin{y}{(f((y)_0) = x \land g((y)_1) = x)}.
\end{align*}
அப்போது $A \cup B$ ஆனது $h$ இன் வரையறைப்புலம்; $A \cap B$
ஆனது~$j$ இன் வரையறைப்புலம்.

\begin{explain}
கணிப்பியல் மொழியில் இங்கு நடப்பது இதுவே: $A$, $B$ இரண்டையும்
எண்ணிடும் செயல்முறைகள் கொடுக்கப்பட்டால், ஏதாவது ஓர் எண்ணீட்டில்
உறுப்பு~$x$ ஐத் தேடுவதன் மூலம் $x$ ஆனது $A \cup B$ இல் உள்ளதா
என்று அரைத் தீர்மானிக்கலாம்; இரு எண்ணீடுகளிலும் ஒரே நேரத்தில்
$x$ ஐத் தேடுவதன் மூலம் உறுப்பு~$x$ ஆனது $A \cap B$ இல் உள்ளதா என்று
அரைத் தீர்மானிக்கலாம்.
\end{explain}

இரண்டாவது நிரூபணத்திற்கு, மீண்டும் $A$ ஆனது~$f$ ஆலும் $B$ ஆனது~$g$
ஆலும் எண்ணிடப்படுகின்றன எனக் கொள்க. பின்வருமாறு எடுத்துக்கொள்க:
\[
k(x) = \begin{cases}
f(x/2) & \text{$x$ இரட்டை எண் எனில்} \\
g((x-1)/2) & \text{$x$ ஒற்றை எண் எனில்.}
\end{cases}
\]
அப்போது $k$ ஆனது $A \cup B$ ஐ எண்ணிடுகிறது; $f$, $g$ வழங்கும்
எண்ணீடுகளுக்கு இடையில் $k$ மாறிமாறிச் செல்கிறது என்பதே இதன்
கருத்து. $A \cap B$ ஐ எண்ணிடுவது சற்றுக் கடினம்.
$A \cap B$ வெற்றுக்கணம் எனில் அது எளிதாகவே கணிக்கத்தக்கவாறு
எண்ணிடத்தக்கது. இல்லையெனில் $c$ ஆனது $A \cap B$ இன் ஏதாவது
ஓர் உறுப்பு ஆகட்டும்; $l$ ஐப் பின்வருமாறு வரையறுக்க:
\[
l(x) = \begin{cases}
f((x)_0) & \text{$f((x)_0) = g((x)_1)$ எனில்} \\
c & \text{இல்லையெனில்.}
\end{cases}
\]
கணிப்பியல் மொழியில், $f$, $g$ இன் எண்ணீடுகளிலுள்ள உறுப்புச்
சோடிகளூடாக~$l$ சென்று, அது காணும் ஒவ்வொரு பொருத்தத்தையும்
வெளியிடுகிறது; இல்லையெனில் $c$ ஐ வெளியிட்டுத் தேங்கி நிற்கிறது.

கடைசி நிரூபணத்திற்கு, $A$ ஆனது பகுதிச் சார்பு $m(x)$ இன்
\emph{வரையறைப்புலம்} என்றும், $B$ ஆனது பகுதிச் சார்பு $n(x)$ இன்
வரையறைப்புலம் என்றும் கொள்க. அப்போது $A \cap B$ ஆனது பகுதிச்
சார்பு $m(x) + n(x)$ இன் வரையறைப்புலம்.

\begin{explain}
கணிப்பியல் மொழியில், $m$ நிற்கும் மதிப்புகளின் கணம்~$A$ என்றும்
$n$ நிற்கும் மதிப்புகளின் கணம்~$B$ என்றும் இருந்தால், இரு
செயல்முறைகளும் நிற்கும் மதிப்புகளின் கணமே $A \cap B$.
\end{explain}

$A \cup B$ ஐ நிற்கும் மதிப்புகளின் கணமாக வெளிப்படுத்துவது இன்னும்
கடினம்; ஏனெனில் $m$ ஐயும் $n$ ஐயும் இணையாக உருவகப்படுத்த வேண்டும்.
$d$ ஆனது $m$ இன் ஒரு சுட்டெண்; $e$ ஆனது $n$ இன் ஒரு சுட்டெண்
ஆகட்டும்; வேறுவிதமாகச் சொன்னால், $m = \cfind{d}$ மற்றும்
$n = \cfind{e}$. அப்போது பின்வரும் சார்பின் வரையறைப்புலமே
$A \cup B$:
\[
p(x) = \umin{y}{(T(d,x,y) \lor T(e,x,y))}.
\]
\begin{explain}
கணிப்பியல் மொழியில், உள்ளீடு $x$ இல், $m$ அல்லது $n$ இன்
நிற்கும் கணிப்பு ஒன்றை~$p$ தேடுகிறது; அவற்றில் ஒன்றைக் கண்டால்
நிற்கிறது.
\end{explain}
\end{proof}

\end{document}
